\documentclass[11pt,a4paper]{article}

% ── packages ─────────────────────────────────────────────────────────────────
\usepackage[top=2.2cm,bottom=2.5cm,left=2.8cm,right=2.8cm]{geometry}
\usepackage[utf8]{inputenc}
\usepackage[T1]{fontenc}
%\usepackage{microtype}
\usepackage{xcolor}
\usepackage{graphicx}
\usepackage{tcolorbox}
\usepackage{titlesec}
\usepackage{enumitem}
\usepackage{parskip}
\usepackage{booktabs}
\usepackage{hyperref}
\usepackage{lastpage}
\usepackage{fancyhdr}
\usepackage[hang,flushmargin]{footmisc}

\tcbuselibrary{skins,breakable}

% ── colour palette ────────────────────────────────────────────────────────────
\definecolor{NavyBlue}{HTML}{1A4E7A}
\definecolor{Amber}{HTML}{F2A60D}
\definecolor{MidBlue}{HTML}{2E7BB8}
\definecolor{GridGray}{HTML}{E8EDF2}
\definecolor{LabelDark}{HTML}{2C3E50}
\definecolor{LightAmber}{HTML}{FFF8E1}
\definecolor{LightBlue}{HTML}{EAF0F6}
\definecolor{BodyText}{HTML}{2C3E50}
\definecolor{SubText}{HTML}{5D6D7E}

% ── hyperref ─────────────────────────────────────────────────────────────────
\hypersetup{
  colorlinks=true,
  urlcolor=MidBlue,
  linkcolor=NavyBlue,
  citecolor=NavyBlue,
  pdfauthor={Boldwin Mweemba},
  pdftitle={A Cleaner Balance Sheet, or Lasting Capacity? Zambia's Copper Test},
}

% ── section headings with coloured rule ──────────────────────────────────────
\titleformat{\section}[block]
  {\large\bfseries\color{NavyBlue}}
  {}{0pt}
  {}
  [\vspace{2pt}\color{Amber}\hrule height 2pt\vspace{4pt}]

\titlespacing{\section}{0pt}{18pt}{8pt}

\titleformat{\subsection}[block]
  {\normalsize\bfseries\color{LabelDark}}
  {}{0pt}{}
\titlespacing{\subsection}{0pt}{12pt}{4pt}

% ── body text ─────────────────────────────────────────────────────────────────
\color{BodyText}
\setlength{\parskip}{7pt}
\setlength{\parindent}{0pt}
\renewcommand{\baselinestretch}{1.18}

% ── header / footer ──────────────────────────────────────────────────────────
\pagestyle{fancy}
\fancyhf{}
\renewcommand{\headrulewidth}{0pt}
\renewcommand{\footrulewidth}{0.4pt}
\fancyfoot[C]{\small\color{SubText}%
  Boldwin Mweemba \textbullet{} Policy Brief \textbullet{} September 2026
  \hfill Page \thepage\ of \pageref{LastPage}}

% ── tcolorbox styles ─────────────────────────────────────────────────────────
\tcbset{
  keyfindingsbox/.style={
    enhanced,
    breakable,
    colback=LightBlue,
    colframe=NavyBlue,
    arc=4pt,
    left=10pt, right=10pt, top=8pt, bottom=8pt,
    boxrule=1.4pt,
    fonttitle=\bfseries\color{white},
    title={Summary},
    attach boxed title to top left={yshift=-2mm, xshift=6mm},
    boxed title style={colback=NavyBlue, arc=3pt},
    coltitle=white,
  },
  calloutbox/.style={
    enhanced,
    colback=LightAmber,
    colframe=Amber,
    arc=4pt,
    left=8pt, right=8pt, top=6pt, bottom=6pt,
    boxrule=1pt,
  },
}

% ─────────────────────────────────────────────────────────────────────────────
\begin{document}
\thispagestyle{fancy}

% ── TITLE BLOCK ──────────────────────────────────────────────────────────────
\begin{center}

{\fontsize{22}{27}\selectfont\bfseries\color{NavyBlue}
A Cleaner Balance Sheet, or Lasting Capacity?\\[4pt]
Zambia's Copper Test}

\vspace{10pt}
\color{Amber}\hrule height 2.5pt
\vspace{10pt}

{\large\color{LabelDark}
\textbf{Boldwin Mweemba}}\\[3pt]
{\normalsize\color{SubText}
Independent Analyst, Data Analytics and Economic Research, Lusaka, Zambia\\
\href{mailto:boldwin@aims.edu.gh}{boldwin@aims.edu.gh}
\quad\textbullet\quad September 2026\\[4pt]
{\small\color{SubText}\textit{This piece consolidates a set of independent
analytical projects on Zambia's economy available at
\href{https://github.com/BoldwinMax}{github.com/BoldwinMax}.}}}

\end{center}

\vspace{14pt}

% ── KEY FINDINGS BOX ─────────────────────────────────────────────────────────
\begin{tcolorbox}[keyfindingsbox]
\begin{itemize}[leftmargin=14pt, itemsep=4pt, topsep=2pt]
  \item Copper has broken \$13,000 per tonne. The Kwacha has appreciated
        32\% since December 2024. Mineral rents have climbed from 3\% to
        nearly 28\% of GDP. The recovery is real and the data confirms it.
  \item The 2020 Eurobond default and its June 2024 restructuring collapsed
        debt service to exports from a 12.5\% peak to under 3\%. That
        improvement is structural, not just cyclical.
  \item Energy is the constraint copper prices cannot relax: the mining
        sector consumed 10.6\% less electricity in 2024 alone, during a
        hydrological crisis that pushed Lake Kariba to 10\% capacity.
  \item Revenue is concentrated dangerously: large taxpayers carry 79\% of
        the domestic tax base; one mine (Kansanshi) contributed roughly a
        third of the entire mining sector's government payments in 2023.
  \item The window to convert this boom into lasting capacity is now, while
        the tailwind is still blowing. Energy, tax diversification, and
        export logistics are the three policy levers that will decide the
        outcome.
\end{itemize}
\end{tcolorbox}

\vspace{10pt}

% ── INTRODUCTION ─────────────────────────────────────────────────────────────
\section{Introduction}

Zambia's economy has been on a run. Copper broke \$13,000 a tonne. The
Kwacha appreciated 32\% between December 2024 and May 2026. Mineral rents
climbed from 3\% of GDP in 2020 to nearly 28\% by 2024. Foreign direct
investment surged past 9\% of GDP once the Eurobond restructuring closed.
Debt service to exports, which peaked at 12.5\% in 2019, has fallen under 3\%.

By almost every headline number, the recovery is real. The question this
piece asks is not whether copper is delivering; the data confirms it is.
The question is whether Zambia is converting that tailwind into lasting
capacity, or simply building a cleaner balance sheet ahead of the next
stress test.

Three things will decide which one it turns out to be: energy reliability,
tax predictability, and export logistics. Copper gives Zambia the
opportunity. Policy determines whether the opportunity becomes
transformation.

% ── FIGURE 1 ─────────────────────────────────────────────────────────────────
\begin{center}
\includegraphics[width=\linewidth]{fig1_transmission.png}
\end{center}
{\small\color{SubText}
\textbf{Figure 1.} The copper to currency transmission chain. Export
earnings growth drives FX reserve accumulation, which in turn supports
Kwacha appreciation and expands fiscal revenue. Key figures at each stage:
\$13,000/tonne copper price; Kwacha appreciation of 32\% (December 2024 to
May 2026); mineral rents rising from 3\% to 28\% of GDP.}

\vspace{6pt}

% ── RESULTS ──────────────────────────────────────────────────────────────────
\section{What the Data Confirms}

The transmission chain from copper to currency is observable and consistent
with the mechanism you would expect: as copper prices rose past
\$13,000 per tonne, export earnings followed, reserves built, and the
Kwacha appreciated in the same window. Foreign exchange reserves alone
explain roughly half of the Kwacha's movements historically, which is more
than copper prices do directly and is itself a useful reminder that the
currency's stability runs through the central bank's reserve position, not
commodity prices alone.

The recovery from the 2020 Eurobond default tells a similar two stage
story. Copper prices nearly doubled between 2020 and 2022, and mineral
rents rose from 3\% to 28\% of GDP before the formal debt restructuring
was finalised. The G20 Common Framework and the June 2024 Eurobond
restructuring then locked in the second stage, collapsing debt service
to exports from 12.5\% to under 3\%. Neither stage tells the full story
alone: without copper, the debt deal would have landed on harsher terms;
without the deal, the copper windfall would have been absorbed by debt
service instead of building reserves.

% ── ENERGY SECTION ───────────────────────────────────────────────────────────
\section{The Constraint the Headline Numbers Don't Show}

Copper mining consumes close to 69\% of Zambia's national electricity.
Reservoir levels at Lake Kariba explain roughly 53\% of the variation in
load shedding hours, and when the reservoir drops, mine electricity
consumption does not hold steady. It fell 10.6\% in 2024 alone, during the
hydrological crisis.

Zambia has set a production target of 3 million tonnes of copper. Getting
there requires a level of consistent power supply the current generation
mix (still roughly 85\% hydro) has not reliably delivered in drought
years. Solar capacity has grown fast, from 88 MW in 2021 to roughly
841 MW now, and that growth has genuinely reduced Zambia's exposure to a
poor rainy season. It has not removed that exposure. Solar does not cover
the evening peak without storage, and hydro still carries most of the load.

The uncomfortable version of this: a country can have record copper prices
and still fail to capture them fully, if the power to run the mines is not
there when it is needed.

\begin{figure}[ht]
\centering
\includegraphics[width=\linewidth]{fig2_energy_v2.png}
{\small\color{SubText}
\textbf{Figure 2.} Energy is the binding constraint copper prices do not
relax. \textbf{Panel A}: Lake Kariba reservoir level (\% full), 2020 to
2024, showing the drop to 10\% during the 2024 hydrological crisis.
\textbf{Panel B}: Load shedding hours per day (red bars) and mine
electricity consumption index (2020\,=\,100), showing the 10.6\% drop in
mine power use in 2024 even as copper prices climbed.}
\end{figure}

% ── TAX BASE SECTION ─────────────────────────────────────────────────────────
\section{The Tax Base Carrying the Weight}

The Zambia Revenue Authority is expected to fund up to 80\% of the national
budget by 2028. Right now, that revenue rests on a narrow base. Large
taxpayers alone carry 79\% of domestic tax revenue. Private businesses
carry 92\% of it. Mining is the largest single contributor, and it also
received 90\% of every tax incentive approved between 2018 and 2023. In
2023 alone, for every kwacha mining paid in tax, it received roughly
8 cents back in incentives. One mine, Kansanshi, paid ZMW 6.3 billion in
2023, which is close to a third of the entire sector's government
contribution on its own.

That concentration is efficient to administer. It is not resilient. A
revenue base this narrow moves with the fortunes of a handful of firms and
a single commodity cycle, which is a fragile foundation for a target as
large as funding 80\% of the national budget.

\begin{figure}[ht]
\centering
\includegraphics[width=0.85\linewidth]{fig3_tax_v2.png}
{\small\color{SubText}
\textbf{Figure 3.} Tax revenue is concentrated in a handful of firms and
one commodity. Large taxpayers carry 79\% of the domestic tax base.
Private businesses carry 92\%. The Kansanshi callout illustrates how thin
the diversification is even within the mining sector itself.}
\end{figure}

% ── DISCUSSION ───────────────────────────────────────────────────────────────
\section{Opportunity Is Not the Same as Transformation}

None of this is a case for pessimism. Copper is real, and right now it is
working in Zambia's favour. The reserves have built. The debt terms have
improved. The recovery is measurable, not narrative.

But a boom converts into lasting capacity only if the constraints around it
get addressed while the window is open, not after prices turn. Energy
reliability determines whether the 3 million tonne production target is
achievable or capped regardless of price. Tax base diversification
determines whether the ZRA's 80\% by 2028 ambition rests on a foundation
that can actually hold that weight. Export logistics determine whether the
earnings copper generates can actually move to market at the pace the boom
demands.

% ── CONCLUSION ───────────────────────────────────────────────────────────────
\section{Conclusion}

Copper bought Zambia the opportunity. The question is what Zambia does
with it, and the answer is almost entirely a function of decisions made
in the next two to three years, while the commodity tailwind is still
blowing and the fiscal space it has created still exists. The political
economy is harder than the economics: tax incentives are easier to grant
than to renegotiate, grid infrastructure competes with other budget
priorities, and logistics corridors require regional coordination. But the
window in which these choices carry the most weight is the one open right
now, when reserves are building, debt service is low, and the sovereign
has fiscal room to make long term commitments.

The three structural constraints at the centre of this analysis are not
abstract concerns. They are the specific bottlenecks that will determine
whether the 3 million tonne production target is met, whether the ZRA's
80\% budget contribution ambition survives the next commodity softening,
and whether the earnings copper generates actually stay in Zambia. The
copper price environment is favourable in ways that will not persist
indefinitely. Whether the decisions being made right now still look right
when the tailwind fades is the question this piece is putting on the table.

% ── REFERENCES ───────────────────────────────────────────────────────────────
\section{References}

\begin{enumerate}[leftmargin=16pt, itemsep=3pt, label={[\arabic*]}]

\item London Metal Exchange (LME). \textit{Copper Official Cash Settlement
      Price}. LME, 2024 to 2026.

\item Bank of Zambia (BoZ). \textit{Monthly Exchange Rate Statistics:
      Zambian Kwacha to US Dollar}. BoZ, Lusaka, 2024 to 2026.

\item World Bank. \textit{World Development Indicators: Mineral Rents,
      FDI, and GDP Series for Zambia}. Washington DC, 2020 to 2024.

\item International Monetary Fund (IMF). \textit{Zambia: 2024 Article IV
      Consultation and Request for Extended Credit Facility}.
      IMF Country Report, Washington DC, 2024.

\item G20 Common Framework for Debt Treatments. \textit{Zambia Debt
      Treatment Agreement}. Official Communiqué, June 2024.

\item Zambia Revenue Authority (ZRA). \textit{Annual Report and Tax
      Expenditure Statistics}. ZRA, Lusaka, 2021 to 2023.

\item Energy Regulation Board (ERB), ZESCO, and Zambezi River Authority.
      \textit{Electricity Generation, Load Shedding, and Kariba Reservoir
      Level Data}. Lusaka and Harare, 2020 to 2024.

\end{enumerate}

\end{document}
